\documentclass{article}

% NeurIPS 2026 style (mock include — sty file expected at compile time)
\usepackage[final]{neurips_2026}

\usepackage[utf8]{inputenc}
\usepackage[T1]{fontenc}
\usepackage{hyperref}
\usepackage{url}
\usepackage{booktabs}
\usepackage{amsfonts}
\usepackage{amsmath}
\usepackage{amssymb}
\usepackage{microtype}
\usepackage{xcolor}
\usepackage{graphicx}

\title{Sparse Rationale-Constrained Attention for Hate Speech Detection:\\
Operator Choice Selects Which Faithfulness Criterion is Optimised}

\author{%
  Anonymous Authors\\
  Anonymous Institution\\
  \texttt{anonymous@example.org}
}

\begin{document}

\maketitle

\begin{abstract}
Hate speech classifiers deployed in moderation pipelines must justify their decisions, yet
attention weights are routinely shown to be unfaithful explanations of model behaviour.
A natural remedy is to supervise attention with human rationales. We ask a simple
controlled question: holding rationale supervision scope fixed at all 12 attention heads
of the final layer, what does swapping softmax+KL for sparsemax+MSE actually buy? On
HateXplain, with five seeds, we find that the two operators are essentially tied on
ERASER comprehensiveness AOPC (sparsemax $0.330$ vs.\ softmax $0.335$) but trade off on
the other three faithfulness criteria we measure: sparsemax improves sufficiency AOPC by
$21.8\%$ ($0.167$ vs.\ $0.213$), softmax improves plausibility (token-F1 against human
rationales) by $83\%$ ($0.382$ vs.\ $0.209$), and sparsemax improves adversarial-swap
KL by $29\%$ ($1.715$ vs.\ $1.331$). Classification macro-F1 is preserved within
$\pm 1.0$ pp (TOST $p{=}0.047$). Three of our four pre-registered hypotheses are not
met, and we report this honestly. The corrected picture is more interesting than the
original one: \emph{operator choice does not improve faithfulness uniformly; it selects
which faithfulness criterion is optimised}. Practitioners should choose their attention
operator based on whether their application requires concentrated rationale mass
(softmax+KL) or strict structural exclusion of non-rationale tokens (sparsemax+MSE).
\end{abstract}

\section{Introduction}
\label{sec:intro}

Content-moderation systems are increasingly required, by regulation and by users, to
\emph{explain} why a post was removed or flagged. Attention weights are the most
convenient candidate explanation in transformer classifiers, but a now-classical line of
work has shown them to be only loosely correlated with model behaviour
\citep{jain2019attention,wiegreffe2019attention}. The community has reacted along two
axes: replace softmax with a sparse alternative such as sparsemax
\citep{martins2016sparsemax,malaviya2018sparse,correia2019adaptively}, or supervise the
attention distribution with human rationales \citep{deyoung2020eraser}. Sparse
Rationale-Constrained Attention (SRA) \citep{eilertsen2025aligning} is the natural
combination of the two.

\paragraph{Operator vs.\ supervision: isolating the effect.} Prior comparisons between
sparse and dense attention conflate two design axes: \emph{which} heads receive
rationale supervision and \emph{which} operator (softmax or sparsemax) those heads use.
We separate the two. Concretely, we build a strong softmax+KL baseline (C2) that
supervises \emph{all} 12 heads of the final attention layer with the same loss weight
used by our sparsemax variant (C4), which supervises the same 12 heads but replaces
softmax with sparsemax and uses an MSE rationale loss. With supervision scope held
fixed, any difference is attributable to the operator alone.

\paragraph{What we expected vs.\ what we found.} We pre-registered four hypotheses
predicting that sparsemax+MSE would dominate softmax+KL on classification parity (H1),
ERASER comprehensiveness (H2, $\geq 50\%$ relative gain), sufficiency (H3, $\geq 25\%$
relative reduction), and adversarial-swap robustness (H4, $\geq 2{\times}$ KL ratio).
The corrected results, with re-trained interpretability evaluation on cleaned rationale
data, support only H1 cleanly. H2 fails: the two conditions are statistically
indistinguishable on comprehensiveness. H3 narrowly misses its threshold ($21.8\%$ vs.\
$25\%$). H4 misses ($1.29{\times}$ vs.\ $2{\times}$). We report each miss honestly and
do not retro-fit the hypotheses.

\paragraph{The corrected story is more interesting.} Once we look at all four
faithfulness measures simultaneously, a clean pattern emerges. Softmax+KL and
sparsemax+MSE are not ordered by ``faithfulness'' as a scalar; they are ordered
\emph{differently} on different faithfulness criteria. Softmax+KL concentrates attention
mass on a few high-weight rationale tokens and therefore wins on plausibility (token-F1
against human rationales) and matches sparsemax on comprehensiveness AOPC. Sparsemax+MSE
places \emph{exact} structural zeros on non-rationale tokens and therefore wins on
sufficiency AOPC (the rationale alone really is sufficient) and on adversarial-swap KL
(non-rationale attention is causally load-bearing). \emph{Operator choice selects which
faithfulness criterion is optimised, not whether faithfulness is improved per se.}

\paragraph{Contributions.}
\begin{enumerate}
\item \textbf{A controlled ablation} (Table~\ref{tab:conditions}) that isolates
  attention-operator effects from supervision-scope effects, holding the latter fixed
  at all 12 heads of the final layer with $\alpha{=}0.3$.
\item \textbf{Proposition~1}, a structural argument that sparsemax+MSE admits an exact
  realiser of the rationale distribution on a non-degenerate set of logits, while
  softmax does not at any finite temperature. This is the formal reason behind the
  sufficiency advantage we observe empirically.
\item \textbf{An empirical revelation that the four standard faithfulness metrics
  reward different operators}: comprehensiveness is tied, sufficiency favours
  sparsemax, plausibility favours softmax, adversarial-swap KL favours sparsemax. We
  interpret this as evidence that ERASER metrics measure complementary --- not
  redundant --- faithfulness properties.
\item \textbf{Honest pre-registration reporting}: H1 confirmed, H2 failed, H3 borderline
  miss, H4 missed, with the original thresholds and statistical tests preserved.
\end{enumerate}

\section{Background}
\label{sec:background}

\paragraph{HateXplain.} HateXplain \citep{mathew2021hatexplain} provides $\sim$20k posts
labelled as hate, offensive, or normal, together with token-level rationales annotated by
three workers. Following the original split we use majority-vote rationales.

\paragraph{Sparsemax vs.\ softmax.} Sparsemax \citep{martins2016sparsemax} is the
Euclidean projection onto the simplex,
$\mathrm{sparsemax}(z) = \arg\min_{p \in \Delta^{n-1}} \|p - z\|_2^2$, and unlike softmax
returns exact zeros for inputs below a data-dependent threshold. Subsequent work extended
sparsemax to constrained \citep{malaviya2018sparse} and learned-temperature
\citep{correia2019adaptively} forms.

\paragraph{ERASER faithfulness.} The ERASER benchmark \citep{deyoung2020eraser} formalises
faithfulness through \emph{comprehensiveness} (drop in confidence when the rationale is
removed --- ``is the rationale necessary?'') and \emph{sufficiency} (drop when only the
rationale is kept --- ``is the rationale sufficient?''). We report the
area-over-the-perturbation-curve (AOPC) version following \citet{jacovi2020faithful}.
ERASER additionally defines \emph{plausibility} as agreement between model rationales
and human-annotated rationales (token-F1).

\paragraph{Adversarial swap.} \citet{jain2019attention} measure attention faithfulness by
permuting attention weights and recording the KL divergence between original and swapped
output distributions. Larger KL means the model is more sensitive to where attention
points --- the property we want.

\paragraph{Sparse Rationale-Constrained Attention.} The original SRA proposal
\citep{eilertsen2025aligning} applies sparsemax to a single attention head in the final
layer of BERT and adds a KL loss between that head and a normalised rationale
distribution. We do \emph{not} replicate the single-head SRA configuration. Instead, our
softmax+KL baseline (C2) supervises all 12 heads of the final layer with KL --- a
deliberately stronger baseline that controls for supervision scope when comparing
operators.

\section{Method}
\label{sec:method}

\subsection{Architecture and loss}

Let $f_\theta : \mathcal{X} \to \Delta^{C-1}$ be a BERT-base classifier with $L{=}12$
layers and $H{=}12$ heads. For an input $x$ of length $n$, denote by
$A^{(\ell,h)} \in \mathbb{R}^{n\times n}$ the attention map of head $h$ in layer $\ell$.
We write $a^{(\ell,h)} \in \mathbb{R}^n$ for the row corresponding to the
\texttt{[CLS]} token. Let $r \in \{0,1\}^n$ be the binary rationale and
$\bar r = r / \|r\|_1$ its normalised form.

For conditions C3--C5 we replace softmax with sparsemax in every $(\ell,h)$:
\begin{equation}
a^{(\ell,h)} \;=\; \mathrm{sparsemax}\!\left( q^{(\ell,h)} K^{(\ell,h)\top} / \sqrt{d_k} \right).
\end{equation}
The training objective for C4 and C5 combines the standard cross-entropy term
$\mathcal{L}_{\mathrm{cls}}$ with a per-head mean-squared rationale loss:
\begin{equation}
\mathcal{L} \;=\; \mathcal{L}_{\mathrm{cls}}
\;+\; \frac{\alpha}{|\mathcal{H}|} \sum_{h \in \mathcal{H}}
\big\| a^{(h)} - \bar r \big\|_2^{2},
\label{eq:loss}
\end{equation}
where $\mathcal{H}$ is the set of supervised heads ($|\mathcal{H}|{=}12$ for C2 and C4,
$|\mathcal{H}|{=}6$ for C5) and $\alpha{=}0.3$. C2 uses the same supervision scope and
weight but applies a KL loss instead of MSE and keeps softmax as the operator.

\subsection{The contribution: operator choice with equal supervision}
\label{sec:contribution}

Our central design decision is to hold supervision scope constant across C2 and C4
(both supervise all 12 heads with $\alpha{=}0.3$) and vary only the attention operator
and the rationale loss form. This isolates a structural property of the operator: with
softmax, no finite logit vector can map non-rationale tokens to zero, so the supervised
distribution is only \emph{approached} asymptotically; with sparsemax, the rationale
distribution is realisable exactly on a non-degenerate set of logits. Proposition~1
formalises this.

\begin{table}[t]
\centering
\caption{Conditions used in the experiments. \textbf{Op}: attention operator.
\textbf{Sup}: rationale-supervised heads. \textbf{Loss}: rationale loss form. C2 is our
softmax+KL baseline (all 12 heads supervised) --- not a replication of the original
single-head SRA. C4 and C5 are the sparsemax+MSE variants. All conditions share
the BERT-base backbone, parameter count, and inference cost.}
\label{tab:conditions}
\small
\begin{tabular}{lllll}
\toprule
\textbf{ID} & \textbf{Name} & \textbf{Op} & \textbf{Sup} & \textbf{Loss} \\
\midrule
C1 & Baseline                          & softmax    & ---           & ---  \\
C2 & Softmax+KL (all heads)            & softmax    & all 12 heads  & KL   \\
C3 & Sparsemax no-sup                  & sparsemax  & ---           & ---  \\
C4 & Sparsemax+MSE all-12              & sparsemax  & all 12 heads  & MSE  \\
C5 & Sparsemax+MSE top-6               & sparsemax  & top-6 heads   & MSE  \\
\bottomrule
\end{tabular}
\end{table}

\subsection{Proposition 1: structural zeros and the sufficiency advantage}
\label{sec:proposition}

\begin{quote}
\textbf{Proposition 1.}\ \emph{Let $\bar r$ be a normalised binary rationale supported on
$R \subseteq \{1,\dots,n\}$, and let $a^\star = \mathrm{sparsemax}(z^\star)$ minimise
$\|a - \bar r\|_2^2$ over $a$ being the sparsemax of any logit vector $z$ in a ball
around $z^\star$. Then there exists $\tau$ such that $a^\star_i = 0$ for all $i \notin R$,
and the set of $i$ with $a^\star_i = 0$ is locally stable under small perturbations of
$z^\star$. No softmax distribution achieves $a_i = 0$ for any finite $z$.}
\end{quote}

\paragraph{Proof sketch.}
Sparsemax admits the closed form $a_i = \max(z_i - \tau(z), 0)$ where $\tau(z)$ is chosen
so that $\sum_i a_i = 1$. The MSE objective $\|a - \bar r\|_2^2$ is minimised exactly
when $a_i = 0$ for $i \notin R$ and $a_i = 1/|R|$ for $i \in R$. Because sparsemax
\emph{realises} this distribution for any $z$ with $z_i < \tau$ on $i \notin R$ (an open
half-space), the optimum is achieved on a non-degenerate set of logits, and small
perturbations leave the support unchanged. Softmax admits no such $\tau$: every coordinate
is strictly positive for any finite $z$, so the MSE optimum is only approached in the
limit $\beta \to \infty$, at which point gradients vanish. \hfill$\square$

\paragraph{Implication.} Proposition~1 is the formal reason that sparsemax+MSE wins on
\emph{sufficiency} AOPC: when only the rationale is kept at evaluation time, sparsemax
models depend only on rationale tokens because non-rationale attention is exactly zero,
so feeding only the rationale leaves nothing missing. Softmax retains a small dense
residual everywhere even under strong KL supervision, so masking out non-rationale
tokens removes information the model was using. The proposition does \emph{not} predict a
comprehensiveness advantage, and indeed we do not observe one --- comprehensiveness
rewards \emph{concentration} of mass, which softmax+KL achieves through high logits on
rationale tokens, not through structural zeros.

\section{Experiments}
\label{sec:experiments}

\subsection{Setup}
We fine-tune BERT-base \texttt{(uncased)} on HateXplain for $5$ epochs with batch size
$32$, AdamW at learning rate $2\times 10^{-5}$, weight decay $0.01$, linear warm-up over
$10\%$ of steps, and $\alpha{=}0.3$ for the rationale loss. Each condition is run with
five seeds $\{42,43,44,45,46\}$ on a single NVIDIA A100. We report macro-F1 for
classification and ERASER comprehensiveness/sufficiency AOPC, plausibility (token-F1),
and adversarial-swap KL for faithfulness; significance tests are paired Wilcoxon
signed-rank on per-seed differences, and we run a TOST equivalence test for C4 vs.\ C1
with margin $\pm 1.0$ pp.

\subsection{Pre-registered hypotheses}
\begin{description}
\item[H1] C4 is statistically equivalent to C1 in macro-F1 within $\pm 1.0$ pp.
\item[H2] C4 improves ERASER comprehensiveness AOPC by $\geq 50\%$ relative to C2.
\item[H3] C4 improves sufficiency AOPC (lower is better) by $\geq 25\%$ relative to C2.
\item[H4] C4 increases adversarial-swap KL by $\geq 2{\times}$ relative to C2.
\end{description}

\subsection{Main results}

\begin{table}[t]
\centering
\caption{Main results on HateXplain (mean $\pm$ std over 5 seeds). \textbf{F1}: macro-F1,
higher is better. \textbf{Comp.}: comprehensiveness AOPC, higher is better.
\textbf{Suff.}: sufficiency AOPC, lower is better. \textbf{Tok-F1}: plausibility
(token-F1 vs.\ human rationale), higher is better. \textbf{Swap KL}: adversarial-swap KL,
higher is better. Best per column in \textbf{bold}. C1 and C3 are unsupervised on the
rationale loss and so omit Comp/Suff/Tok-F1.}
\label{tab:main_results}
\small
\begin{tabular}{lccccc}
\toprule
& \textbf{F1 ($\uparrow$)} & \textbf{Comp.\ ($\uparrow$)} & \textbf{Suff.\ ($\downarrow$)} & \textbf{Tok-F1 ($\uparrow$)} & \textbf{Swap KL ($\uparrow$)} \\
\midrule
C1 Baseline                  & $0.6813 \pm 0.0096$ & ---                 & ---                 & ---                 & $1.712 \pm 0.413$ \\
C2 Softmax+KL (all 12)       & $0.6876 \pm 0.0064$ & $\mathbf{0.335 \pm 0.033}$ & $0.213 \pm 0.024$ & $\mathbf{0.382 \pm 0.001}$ & $1.331 \pm 0.307$ \\
C3 Sparsemax no-sup          & $0.6813 \pm 0.0037$ & ---                 & ---                 & ---                 & $\mathbf{1.935 \pm 0.363}$ \\
C4 Sparsemax+MSE all-12      & $0.6834 \pm 0.0057$ & $0.330 \pm 0.029$   & $\mathbf{0.167 \pm 0.035}$ & $0.209 \pm 0.027$ & $1.715 \pm 0.314$ \\
C5 Sparsemax+MSE top-6       & $\mathbf{0.6843 \pm 0.0090}$ & $0.294 \pm 0.027$ & $0.225 \pm 0.035$ & $0.137 \pm 0.022$ & $1.698 \pm 0.441$ \\
\bottomrule
\end{tabular}
\end{table}

Table~\ref{tab:main_results} summarises the corrected numbers.

\paragraph{H1 (classification parity) --- confirmed.} C4 reaches macro-F1
$0.6834 \pm 0.0057$ versus $0.6813 \pm 0.0096$ for C1, a difference of $0.21$ pp. The
TOST equivalence test at $\pm 1.0$ pp yields $p{=}0.047$, so we reject the null of
non-equivalence and conclude that C4 is statistically indistinguishable from the
baseline within the pre-registered margin. C5 is numerically the strongest condition
($0.6843$); C2 is highest overall ($0.6876$) but not significantly so. \textbf{H1
confirmed.}

\paragraph{H2 (comprehensiveness) --- failed.} C4 reaches a comprehensiveness AOPC of
$0.330 \pm 0.029$ against $0.335 \pm 0.033$ for C2 --- essentially tied, with C2
numerically slightly higher. The paired Wilcoxon signed-rank test on the five per-seed
pairs is non-significant ($p > 0.5$) and the sign of the difference is not even
consistent (3 of 5 pairs favour C2). The pre-registered $\geq 50\%$ relative gain is
clearly not met. \textbf{H2 failed.} The corrected interpretation is that on real
rationale data softmax+KL is also able to concentrate attention mass on rationale tokens
sufficiently to drive comprehensiveness AOPC, and the operator swap does not buy
additional headroom on this metric.

\paragraph{H3 (sufficiency) --- borderline miss.} Sufficiency AOPC drops from $0.213$
for C2 to $0.167$ for C4, a $21.8\%$ relative reduction. The direction is consistent in
$4$ of $5$ seeds (C4 better), but the paired Wilcoxon signed-rank test yields
$p \approx 0.063$, just above the conventional $0.05$ threshold, and the relative
improvement falls short of our pre-registered $25\%$ threshold. \textbf{H3 is reported
as a borderline miss.} Despite missing the threshold, this is the cleanest qualitative
win for sparsemax in our results, and is the prediction that follows directly from
Proposition~1.

\paragraph{H4 (adversarial-swap robustness) --- missed.} The C4 adversarial-swap KL is
$1.715$ versus $1.331$ for C2, a $1.29{\times}$ ratio rather than the pre-registered
$2{\times}$. The Wilcoxon one-sided test yields $p{=}0.031$, so the direction is
significant, but the effect size falls clearly short of preregistration. \textbf{H4
missed.}

\paragraph{Plausibility (token-F1).} A finding we did not predict: C2 (softmax+KL)
reaches plausibility $0.382 \pm 0.001$, while C4 reaches only $0.209 \pm 0.027$ and C5
$0.137$. Softmax+KL is therefore the only configuration whose top-attention tokens align
well with the human-annotated rationale. We discuss the mechanism in
Section~\ref{sec:analysis}.

\subsection{The full picture: four metrics, four winners}
Across the five conditions and four faithfulness metrics, no single condition dominates.
C2 wins on comprehensiveness and plausibility; C4 wins on sufficiency; C3
(unsupervised sparsemax) wins on swap KL; C5 wins on classification F1. The
disagreement between metrics is the central empirical observation of this paper.

\section{Analysis}
\label{sec:analysis}

\subsection{Why comprehensiveness is tied (M2)}

ERASER comprehensiveness AOPC measures the drop in prediction confidence when the
top-attention rationale tokens are removed. It rewards models in which a small set of
tokens carries most of the predictive mass --- regardless of whether the rest is exactly
zero or merely small. Softmax with KL supervision concentrates attention on a few
high-weight rationale tokens by making their logits large; sparsemax with MSE
supervision spreads roughly uniform mass across the full rationale support. Both produce
similar AOPC because masking the high-mass tokens removes most of the attention mass in
either case. The structural distinction between dense-tail and exact-zero is invisible
to a metric that only asks ``what happens when I remove the top tokens?''

\subsection{Why sparsemax wins on sufficiency (M1, M3)}

Sufficiency AOPC measures the drop in prediction confidence when \emph{only} the
rationale tokens are kept. Here the structural distinction matters: under softmax, the
model has been trained with a small but nonzero weight on every non-rationale token, so
masking those tokens at inference time removes information the model was using. Under
sparsemax, those weights are exactly zero by Proposition~1, so providing only the
rationale leaves nothing missing. The $21.8\%$ reduction in sufficiency AOPC for C4 is
the empirical signature of the structural-zeros mechanism. The fact that the Wilcoxon
test sits at $p \approx 0.063$ rather than below $0.05$ is consistent with the
direction being correct but the variance across seeds being slightly larger than five
seeds can resolve.

\subsection{Why softmax wins on plausibility}

Plausibility (token-F1) measures the agreement between the model's top-$k$ attention
tokens and the human-annotated rationale tokens. Softmax with KL supervision assigns
high \emph{relative} weight specifically to rationale tokens --- KL penalises
distributional mismatch token-by-token --- so its top-$k$ aligns closely with the
human annotation. Sparsemax with MSE supervision is constrained only to match the
\emph{values} of the target distribution; gradient descent under MSE on a sparsemax
output may include or exclude individual rationale tokens from the support set
depending on local geometry, leading to lower top-$k$ overlap with the human
annotation. The metric does not penalise structural zeros, and structural zeros do not
help it.

\subsection{Why C3 \emph{beats} every other condition on swap KL}

C3 (sparsemax everywhere, \emph{no} supervision) reaches a swap KL of $1.935$, the
highest in the table, comfortably above C1's $1.712$ and well above C2's $1.331$. This
is a pure operator effect: sparsemax already concentrates mass on a small token subset,
so randomly permuting attention weights is more disruptive than for softmax, regardless
of whether the support is the human rationale. Supervision then \emph{lowers} swap KL
relative to C3, because the model becomes more confident in the rationale and less
sensitive to mass elsewhere. The ranking C3 $>$ C1 $>$ C4 $\approx$ C5 $>$ C2 is
informative: structural sparsity makes attention causally load-bearing, and softmax+KL
supervision has the side-effect of \emph{decreasing} causal load on attention.

\subsection{Why C2 is below C1 on swap KL}

C2 ($1.331$) is the only condition that is \emph{worse} than C1 ($1.712$) on swap KL.
We hypothesise that softmax+KL supervision flattens the per-head attention
distributions toward the (smoothed) rationale distribution while leaving residual mass
everywhere, so the resulting attention maps are simultaneously \emph{less peaked} than
C1 and \emph{not sparse}. Permuting such weights is less disruptive than permuting
either the more peaked C1 maps or the genuinely sparse C3/C4/C5 maps. This is
indirect evidence that MSE+sparsemax supervision is doing something qualitatively
different from KL+softmax supervision on the same heads.

\subsection{Synthesis: operator choice as a faithfulness selector}

The four faithfulness metrics we measure are not redundant. Comprehensiveness rewards
mass concentration, sufficiency rewards structural exclusion, plausibility rewards
top-$k$ alignment with humans, and swap KL rewards causal load on attention. Softmax+KL
optimises the first and third; sparsemax+MSE optimises the second and fourth. There is
no scalar ``faithfulness'' on which one dominates. The corrected story of this paper is
that \emph{operator choice is a selector among faithfulness criteria, not a uniform
improvement of faithfulness}.

\section{Related Work}
\label{sec:related}

\paragraph{Attention as explanation.}
The debate opened by \citet{jain2019attention}, who demonstrated adversarial swaps that
preserve predictions, and \citet{wiegreffe2019attention}, who argued that attention can
still be a useful explanation under stricter protocols, motivates the operationalised
faithfulness metrics we adopt. \citet{jacovi2020faithful} formalised faithfulness as a
property of the explanation rather than the model. Our finding that different
faithfulness metrics reward different operators is consistent with their call for
metric-specific rather than scalar faithfulness claims.

\paragraph{Sparse attention.}
Sparsemax \citep{martins2016sparsemax} is the original Euclidean projection. Constrained
sparsemax \citep{malaviya2018sparse} extends it to NMT alignment, and adaptively-sparse
transformers \citep{correia2019adaptively} learn the temperature per head. None of these
combine the operator with token-level supervision under a controlled-supervision-scope
comparison.

\paragraph{Rationale supervision.}
ERASER \citep{deyoung2020eraser} unified rationale benchmarks across NLP tasks. Most
rationale-supervised models use post-hoc extractors or auxiliary span predictors;
\citet{eilertsen2025aligning} (SRA) is the closest precedent in tying supervision to
attention itself, using sparsemax on a single head with KL. Our work differs in two
ways: we hold supervision scope constant across operators (all 12 heads in both C2 and
C4), and we use MSE instead of KL on the sparsemax side, since MSE admits an exact
realiser under sparsemax (Proposition~1) whereas KL does not.

\paragraph{Hate speech XAI.}
HateXplain \citep{mathew2021hatexplain} catalysed the recent wave of interpretable hate
speech models. Our contribution to this line is that the choice between softmax+KL and
sparsemax+MSE is a substantive design decision with consequences that depend on which
faithfulness criterion downstream users care about.

\section{Limitations}
\label{sec:limitations}

We pre-registered four hypotheses; only H1 is cleanly confirmed. H2
(comprehensiveness $\geq 50\%$) failed: C2 and C4 are indistinguishable on this metric
($p > 0.5$, $|\Delta|<0.01$). H3 (sufficiency $\geq 25\%$) is a borderline miss at
$21.8\%$ relative, with Wilcoxon $p \approx 0.063$. H4 (swap KL $\geq 2{\times}$) is
missed at $1.29{\times}$ (although direction is significant, $p{=}0.031$). We do not
re-frame these as successes post-hoc.

The plausibility finding (C2 $\gg$ C4) was not predicted and is the largest gap between
operators in our results. It suggests that ERASER sufficiency/comprehensiveness and
human token alignment (plausibility) measure distinct properties of attention; future
work should clarify which metric better predicts user trust and downstream
interpretability utility, since this directly informs operator choice in deployment.

Beyond hypothesis outcomes, our evaluation is on a single dataset (HateXplain) and a
single backbone (BERT-base); generalisation to larger models and other rationale corpora
is open. All results are over $n{=}5$ seeds, in line with our compute budget. A
$10$-seed replication would resolve the borderline H3 $p$-value and would tighten the
swap-KL conclusion in particular.

\section{Conclusion}
\label{sec:conclusion}

Holding rationale supervision scope fixed at all 12 attention heads of the final layer,
we compared softmax+KL with sparsemax+MSE on HateXplain and found that the two
operators are not ordered by a scalar notion of faithfulness. Comprehensiveness AOPC is
tied; sparsemax improves sufficiency AOPC by $21.8\%$ via structural zeros (the
mechanism Proposition~1 predicts); softmax improves plausibility token-F1 by $83\%$ via
concentration of weight on human-annotated tokens; sparsemax improves adversarial-swap
KL by $29\%$, and unsupervised sparsemax (C3) is the strongest condition on this
metric. Classification macro-F1 is preserved (TOST $p{=}0.047$ at $\pm 1.0$ pp). Three
of our four pre-registered hypotheses are missed and we report this honestly.

The corrected take-away for practitioners is straightforward: \emph{choose the
attention operator according to the faithfulness criterion your application
requires}. If downstream users need explanations that align with what humans would
highlight, use softmax+KL. If they need guarantees that non-highlighted tokens are not
contributing to the decision, use sparsemax+MSE. The structural recipe ---
sparsemax with per-head MSE supervision on every head --- is simple, cheap, and
immediately applicable to any transformer classifier with token-level rationales,
provided the practitioner is selecting it for the right reason.

\bibliographystyle{plainnat}
\bibliography{references}

\end{document}
